\section{Conclusion}
This paper presents SoCPilot, an LLM-driven framework for application-driven CPU--accelerator SoC customization. SoCPilot coordinates requirement interpretation, application profiling, hardware/software partitioning, accelerator generation, SoC configuration and integration, and latency--area design space exploration. The generated hardware configurations are synthesized and validated on an FPGA board, where complete-application execution verifies accelerator invocation, communication, and functional correctness. Benchmark-specific embedded runtime software is manually implemented using the ESP programming model and is explicitly kept outside the LLM-generated portion of the flow.

The evaluation is organized around three questions: whether the generated systems provide correct end-to-end acceleration over matched CPU-only implementations, whether a specialized SoC workflow is more reliable than open-ended general-purpose agent execution for accelerator generation and integration, and whether SoCPilot adapts CPU--accelerator configurations under different area constraints. The current results demonstrate the potential of tool-grounded LLM orchestration for cross-layer SoC customization while also making the automation boundary explicit. Future work will extend the automated flow toward embedded-software generation and broaden system-level cost models and validation coverage.

